\clearpage
\thispagestyle{empty}
\pdfbookmark[0]{Über die Daltonsche Theorie. Ausgabe und Quellen}{dalton1830-edition}
\begin{center}\Large\textbf{Besprechung von Benzenbergs \textit{Über die Daltonsche Theorie}}\end{center}

Carl Friedrich Gauss’ vollständige Besprechung von Johann Friedrich Benzenbergs Buch in den \textit{Werken}, Band V (Göttingen, 1867), S.~583--591, liegt hier zusammen mit dem vorangestellten Titel Besprechungen fremder Schriften als KI-Übertragung in modernes Deutsch mit KI-erstelltem Satz vor. Gauss untersucht, wie Benzenberg die Daltonsche Hypothese auf barometrische Höhenmessungen anwandte.

Die Besprechung erschien zuerst in den \textit{Göttingischen gelehrten Anzeigen}, 196. Stück, 11. Dezember 1830, S.~1945--1956. Die Übertragung folgt dem Wortlaut, der Zahlenschreibweise und den zehn Tabellen der Werkausgabe.

\medskip\noindent
\href{https://gdz.sub.uni-goettingen.de/id/PPN236006339}{Deutsche Werkausgabe (1867)}\par
\noindent\href{https://gdz.sub.uni-goettingen.de/id/PPN385030444_1830_3}{Erstdruck der Besprechung (1830)}
\clearpage
